% Transcription — fonds Alexandre Grothendieck, shelfmark 48, batch 1 (pages 1-5).
% Established from the facsimile archives/batches/48/batch-01.pdf (a mirror of
% https://grothendieck.umontpellier.fr/48.pdf), per the skill
% transcribe-grothendieck.
% Pages 1, 2 and 4 are the plan the folder is named for: a cover, a chapter
% outline in lettered parts A to C with numbered sections, and a continuation in
% D and E with a short reading list. Pages 3 and 5 are two leaves of a different
% document altogether — a typescript of SGA 4 carrying its own pagination (29)
% and (30), on the points of a topos — which the author has gone over correcting
% the typist's mathematical symbols. They are mathematics and they are in the
% folder, so the theorem they turn on is transcribed; the body of the typescript
% is not re-set here, since it is printed text and not a hand.
% Pass: Opus 5 (claude-opus-5), 2026-08-16 — first pass, unchecked against the pages by a human.
\documentclass[11pt,a4paper]{article}
% Le préambule commun à toutes les transcriptions du fonds Grothendieck.
%
% Il tient en une idée : **l'incertitude se compose, elle ne se résout pas.**
% Une transcription qui lisse un mot douteux a détruit la seule chose pour
% laquelle on la faisait — un lecteur doit pouvoir savoir, à chaque endroit,
% ce qui était sur le feuillet et ce qui a été ajouté. Les six macros qui
% suivent sont donc l'appareil critique, et non de la décoration.
%
% Les mêmes commandes sont reconnues par `scripts/render.mjs`, qui produit la
% vue de lecture du site. Une macro ajoutée ici doit l'être aussi là-bas,
% faute de quoi le rendu échouera bruyamment — ce qui est voulu : un rendu
% silencieusement faux est pire qu'un rendu absent.

\ProvidesPackage{grothendieck}[2026/08/08 Transcriptions du fonds Grothendieck]

\usepackage{amsmath,amssymb,amsthm}
\usepackage{mathrsfs}
\usepackage{stmaryrd}
% Les diagrammes commutatifs sont partout dans ce fonds : c'est la langue
% même de la géométrie algébrique de Grothendieck, et une transcription qui
% les rendrait en prose ne transcrirait rien.
\usepackage{tikz-cd}
% Français par défaut — c'est la langue des feuillets — mais l'anglais est
% chargé aussi : la traduction et le résumé sont des documents anglais, et la
% typographie française (espaces avant les deux-points) y serait une faute.
% Ces documents basculent par \selectlanguage{english} en tête de corps.
\usepackage[english,french]{babel}
\usepackage{fontspec}
\usepackage{geometry}
\geometry{a4paper,margin=25mm,marginparwidth=28mm}
\usepackage{marginnote}
\usepackage{xcolor}
% ulem, not soul. `soul`'s \st and \ul are built on a character-by-character
% scan that XeLaTeX's font machinery breaks: the document fails with an
% undefined control sequence rather than degrading. ulem does the same two jobs
% and is Unicode-engine safe. `normalem` stops it hijacking \emph, which the
% transcriptions use for Grothendieck's own emphasis.
\usepackage[normalem]{ulem}
\usepackage{hyperref}
\hypersetup{colorlinks=true,linkcolor=black,urlcolor=black,pdfborder={0 0 0}}

% --- Métadonnées du lot -----------------------------------------------------
% Reprises telles quelles par le script de rendu, qui les affiche en tête de
% la vue de lecture. Elles fixent ce que le fichier prétend couvrir : sans
% elles, une transcription flotte sans point d'attache dans un fonds de seize
% mille feuillets.

\newcommand{\folder}[1]{\gdef\grothFolder{#1}}
\newcommand{\batch}[1]{\gdef\grothBatch{#1}}
\newcommand{\leaves}[2]{\gdef\grothFirst{#1}\gdef\grothLast{#2}}
\newcommand{\foldertitle}[1]{\gdef\grothFoldertitle{#1}}
\gdef\grothFolder{?}\gdef\grothBatch{?}\gdef\grothFirst{?}\gdef\grothLast{?}\gdef\grothFoldertitle{}

% --- Le résumé --------------------------------------------------------------
%
% Il ouvre la lecture modernisée, sous le titre « Résumé », et il est composé
% en petit. Ce n'est pas un ornement : c'est le seul passage du document qui
% ne soit pas des mathématiques, et un lecteur doit pouvoir le reconnaître
% avant de l'avoir lu — puis le sauter s'il connaît déjà le sujet, ou s'y
% arrêter s'il ne le connaît pas. Le filet qui le suit marque la reprise du
% corps.

\newenvironment{resume}
  {\small\setlength{\parskip}{0.45em}}
  {\par\medskip\hrule\medskip}

% --- Mots-clés ---------------------------------------------------------------
%
% En anglais, en clôture du résumé de la lecture modernisée : le vocabulaire
% moderne sous lequel chercher ce que les feuillets construisent. Le manifeste
% du site (`npm run manifest`) les extrait pour étiqueter la cote entière —
% cette ligne est la source unique des tags, il n'y a pas de fichier à côté.
\newcommand{\keywords}[1]{\par\smallskip\noindent
  {\footnotesize\textsc{Keywords} --- \textit{#1}}\par}

% --- L'appareil critique ----------------------------------------------------

% Le numéro de feuillet, en marge. C'est l'ancre qui permet de régler un
% désaccord sur une formule en regardant *un* feuillet plutôt qu'une chemise
% de six cents. Le site s'en sert aussi pour tourner les pages du fac-similé
% quand on fait défiler la transcription.
\newcommand{\leaf}[1]{%
  \marginnote{\footnotesize\color{gray!70}\textsf{#1}}%
  \label{leaf:#1}\ignorespaces}

% Illisible : on ne devine pas. Un mot inventé qui se lit comme les autres est
% le pire résultat possible.
\newcommand{\ill}{\textcolor{red!60!black}{[\,\ldots\,]}}

% Lecture douteuse : le mot est donné, et signalé comme incertain.
\newcommand{\uncertain}[1]{\uline{#1}}

% Ajout de l'éditeur — une lettre manquante, un mot sous-entendu.
\newcommand{\add}[1]{\textcolor{blue!50!black}{[#1]}}

% Biffé par Grothendieck lui-même. Ce qu'il a rayé fait partie du document :
% une rature dit souvent où la pensée a changé de direction.
\newcommand{\struck}[1]{\sout{#1}}

% Note du transcripteur, sur l'état matériel du feuillet.
\newcommand{\note}[1]{\par\smallskip{\small\color{gray!60!black}\textit{[#1]}}\par\smallskip}

% Note marginale de Grothendieck, distincte du corps du texte.
\newcommand{\marginal}[1]{\par\smallskip\noindent\hspace{1em}%
  {\small\color{gray!60!black}#1}\par\smallskip}

% --- Titre ------------------------------------------------------------------

\AtBeginDocument{%
  \begin{center}
    {\large\bfseries Fonds Alexandre Grothendieck --- cote n\textsuperscript{o}~\grothFolder}\\[2pt]
    {\itshape\grothFoldertitle}\\[4pt]
    {\small feuillets \grothFirst--\grothLast\ (lot \grothBatch)}\\[2pt]
    {\footnotesize\color{gray!60!black}Transcription établie d'après le fac-similé numérisé
     par l'Université de Montpellier.\\
     Les lectures incertaines sont soulignées, les illisibles notées [\,\ldots\,],
     les ajouts entre crochets.}
  \end{center}
  \vspace{1em}}

\endinput


\folder{48}
\batch{1}
\pages{1}{5}
\dating{[à partir de 1967]}
\watermark{Édition de démonstration}
\foldertitle{Plan EGA VII : notes manuscrites (s.d.)}

\begin{document}

\page{1}
En haut à droite, à l'encre, le chiffre romain souligné et surligné :
« Plan EGA VII ». Le reste de la page est vierge.

\page{2}

\section*{EGA Chap.\ VI --- Schémas en groupes et \add{torseurs} \struck{\ill}}

\note{le titre porté ici est « Chap.\ VI », alors que la chemise annonce
« Plan EGA VII » : les deux numéros sont sur les pièces du même dossier et ne
sont pas rendus concordants. Après « et », un mot ou deux ont été rayés au
point d'être illisibles et « torseurs » écrit au-dessus}

\subsection*{(A) Propriétés générales}

\begin{itemize}
\item[§1] Histoire du passage d'un \ill\ à travers, des propriétés ouvertes /
propriétés générales des morphismes

\item[§2] Composantes neutres. Dimensions. Séparation

\item[§3] Groupes affines

\item[§4] Sous-groupes engendrés \add{\ill}

\item[§5] Passage au quotient \ill\ et \ill\ quotients

\item[§6] $G^{\tau}$ et $G$ \ill\ (cf.\ \ill\ VI n\textsuperscript{o} 1)

\note{le sigle abrégé de ce renvoi est écrit en capitales grasses de quatre
lettres ; il n'est pas assuré et n'est pas restitué. Le renvoi du §7 ci-dessous,
lui, se lit SGAD}
\end{itemize}

\marginal{une accolade tracée dans la marge de gauche réunit les §3 à §6, avec
un mot en regard resté illisible}

\subsection*{(B) Étude infinitésimale}

\begin{itemize}
\item[§7] Questions d'extension infinitésimale. (cf.\ SGAD II)

\item[§8] Hyperalgèbre d'un \ill\ et gr.\ \uncertain{eff.} \ill

\item[§9] Prolongements

\item[§10] Algèbres \add{de} groupes réduites et \ill\ de Lie

\item[§11] Dualité de Cartier

\item[\struck{§\ill}] \struck{\ill}
\end{itemize}

\subsection*{(C) Groupes de \ill ; propriétés de rigidité liées aux points d'ordre fini}

\begin{itemize}
\item[§12] Maximum : SGAD VIII, IX, X
\end{itemize}

\page{4}

\subsection*{(D) Propriétés générales des schémas abéliens et polarisations}

\subsection*{(E) Extensions des schémas abéliens par groupes finis (th.\ de \uncertain{\ill}\ldots)}

\note{les deux lettres D et E sont reprises plus bas, une seconde fois et
doublement cerclées, pour deux autres rubriques : la page compte donc deux D et
deux E, et rien n'indique laquelle des deux séries l'emporte}

\subsection*{(D) \struck{Groupes de \ill} Torseurs sous gr.\ classiques (n'importe où)}

et théorème d'irréductibilité \ill\ (P.)

\subsection*{(E) Modèles de Néron (?)}

Un paragraphe de variantes, avec le « th.\ des carrés » \ill

\subsection*{Réf.}

\begin{itemize}
\item SGAD
\item BIBLE
\item « Boulder »
\item Livre(s) Mumford
\end{itemize}

\page{5}

\note{feuillet de tapuscrit, paginé (29) en haut à droite dans sa propre
numérotation, repris à la main : les corrections portent presque toutes sur les
lettres grecques et les accents circonflexes que la machine n'avait pas, et
qui sont rétablis à la plume dans la marge et dans le corps. On donne ici
l'énoncé qui porte le feuillet ; le corps du tapuscrit n'est pas re-composé}

\[
(*) \qquad \underline{\mathrm{Pt}}(X)^{\circ} \longrightarrow
\underline{\mathrm{Pt}}(\widetilde{X_{\text{ét}}})
\]

\textbf{Théorème 7.9.} Soit $X$ un préschéma. Le foncteur précédent $(*)$ est
une anti-équivalence de la catégorie des points géométriques sur $X$ (avec
comme morphismes les flèches des spécialisation) avec la catégorie des
foncteurs fibres sur le topos étale $\widetilde{X_{\text{ét}}}$.

\note{le « anti » de « anti-équivalence » est repassé à la main sur une
frappe fautive}

La démonstration se poursuit en a) \emph{Le foncteur envisagé est pleinement
fidèle}, où l'on pose, pour $X \in \mathrm{Ob}\,X_{\text{ét}}$,
$\check{X} : \widetilde{X_{\text{ét}}} \to (\mathrm{Ens})$ défini par
$\check{X}(F) = \mathrm{Hom}(X, F) = F(X)$, et où le foncteur fibre
$\varepsilon_{\xi}$ s'écrit
$\varepsilon_{\xi} \simeq \varinjlim_{i} \check{X_{i}}$, les $X_{i}$ étant des
préschémas affines étales sur $X$ indexés par une catégorie
pseudo-filtrante.

\page{3}

\note{second feuillet du même tapuscrit, paginé (30), qui continue la
démonstration : la comparaison de $(*)$ et $(**)$, puis b) \emph{Le foncteur
envisagé est essentiellement surjectif}, 7.9.1. Les reprises à la main y sont
de même nature qu'au feuillet précédent --- des lettres grecques rétablies ---
sauf une, portée dans la marge en regard de b) et qui se lit
\uncertain{« \ldots\ d'enfoncement »}}

\note{les deux feuillets sont classés dans cet ordre par l'archiviste, 3 avant
5, alors que leur pagination propre les donne dans l'autre sens : (29) est la
page 5 et (30) la page 3. Ils sont donnés ici dans l'ordre de leur
démonstration}

\end{document}
